\subsection{Exact completion of the quartic centralizer data}
\label{subsec:frontier-quartic-completion}

The full tensor of $R_{6,4,91}$ contains the Ising based subcategory
with $g=b_4$ and $\sigma=b_3$. The preceding signed and induction
reductions leave two possible center ranks, $36$ and $42$. Each profile
has exactly one central simple restricting to $g$ and exactly one
restricting to $\sigma$. Their products, determined by restriction and
simplicity, give an Ising fusion subcategory of the center. Every braided
Ising category is nondegenerate. Proposition~\ref{prop:central-packets-adjunction}
therefore applies to its centralizer $\mathcal E$.

\begin{theorem}[Complete centralizer fusion and modular data for $R_{6,4,91}$]
\label{thm:91-centralizer-completion}
If $R_{6,4,91}$ has a pseudo-unitary categorification $\mathcal C$, then
$Z(\mathcal C)\simeq\mathcal I\boxtimes\mathcal E$, where $\mathcal I$
is Ising and $\mathcal E$ has rank $14$. The center-rank-$36$ profile is
impossible. For the rank-$42$ profile, the complete compatible fusion
ring of $\mathcal E$ has one based-ring isomorphism class. With its
restriction columns labelled as in the certificate, it has eight fusion
tensors, all related by permutations of repeated restriction columns.
There are exactly two compatible pairs of centralizer modular data;
they are complex conjugate. These assertions specify necessary central
data and do not assert that either pair is categorifiable.
\end{theorem}
\begin{proof}
Write the low summands of canonical unit induction as $A_0,\ldots,A_5$.
Their known character eigenvalues at the pure $g$ and $\sigma$ columns
identify their Ising coordinates without choosing a braiding on
$\mathcal C$:
\[
 A_0=1\boxtimes1,\quad A_1=g\boxtimes a,\quad
 A_2=g\boxtimes b,\quad A_3=1\boxtimes c,\quad
 A_4=\sigma\boxtimes u,\quad A_5=\sigma\boxtimes u^*.
\]
Both inherited low character labellings give these same coordinates
and the same duality permutation $(4\ 5)$ on low labels. Thus the right
adjoint of $F|_{\mathcal E}$ at the tensor unit is $1\oplus c$.

The complete central profile must be a union of packets
$(v,gv,\sigma v)$. The nonnegative integral packet-cover equations
have one solution for each of the two profiles. Each admits $144$
endpoint orientations up to permutations of equal columns. Of these,
$135$ violate the displayed low Ising coordinates. For a remaining
orientation, let $V$ have the proposed centralizer restriction columns.
The adjunction identity fixes
\[
 N_c=V^{\mathsf T}V-I.
\]
Nonnegative entries, the unit column, and the tensor-functor equation
$VN_c=N_{F(c)}V$ reject eight further orientations. Exactly one
restriction multiset survives for each rank. This is a decomposition
of the center; no tensor-product decomposition of $\mathcal C$ is used.

We describe the entire finite fusion search. For every possible duality
involution consistent with the restriction columns and the forced pair
$u,u^*$, use the symmetric variables
\[
 x_{ijk}=N_{ij}^{k^*},\qquad x_{ijk}=x_{\pi(i,j,k)}.
\]
The unit and $c$ products are fixed. Restriction of every product gives
linear equations and finite upper bounds
\[
 N_{ij}^{k}\le
 \min_{h:(v_k)_h>0}
 \left\lfloor\frac{(v_iv_j)_h}{(v_k)_h}\right\rfloor.
\]
Duality and commutation with $N_c$ give further exact linear equations.
Project this system onto multiplication by the simple whose restriction
is $b_1$, eliminating the complementary variables by rational row
reduction. Every actual fusion matrix is normal, since multiplication
by its dual is its transpose and the centralizer ring is commutative.

The rank-$12$ centralizer has one possible duality. Its projected affine
system has eight free coordinates and no bounded integral solution;
complete enumeration visits $11$ nodes. For rank $14$, there are four
possible dualities. Three give $64$ projected matrices each, none normal.
The remaining duality gives $512$ projected matrices, of which $16$ are
normal. For each such matrix, its commutation with all the other fusion
matrices is now linear. Eight completions are inconsistent; each of the
other eight has one bounded integral completion. Direct multiplication
checks the unit, duality, nonnegative integral coefficients and every
associativity identity. Swapping the three pairs of identical restriction
columns identifies these eight tensors in one based-ring orbit.

For each orbit representative the exact character table is obtained by
factoring a separating multiplication matrix over $\QQ(\zeta_{16})$.
All character products and all roots are checked exactly. In any modular
labelling, the character attached to object $j$ has formal codegree
$\operatorname{FPdim}(\mathcal E)/d_j^2$, and
$\widehat S_{ij}=d_j\chi_j(i)$ must be symmetric. Exhaustive assignment
of the characters subject to these requirements gives exactly two
$\widehat S$ matrices.

Put $z=\theta_\sigma$. The inherited second indicator is
$\nu_2(\sigma)=-1$, so
$z=\zeta_{16}^{e}$ with $e\in\{3,5,11,13\}$.
The low summands have twist one, hence
\[
 \theta_a=\theta_b=-1,\qquad \theta_c=1,\qquad
 \theta_u=\theta_{u^*}=z^{-1}.
\]
For each $S$ and each of these four values of $z$, apply balancing
with one of these known-twist objects in the first coordinate:
\[
 \theta_i\theta_j\widehat S_{ij}
   =\sum_k N_{i^*,j}^{k}d_k\theta_k.
\]
These are linear equations for all fourteen twists. Every branch has
rank fourteen. Exactly two are consistent, with Ising exponents $5$
and $11$, respectively. The unique solutions have twists in
$\mu_{16}$ and are complex conjugate. In particular, no a priori bound
on the unknown twist orders is being substituted for an exhaustive
search: the linear equations determine the twists themselves.
Each solution passes every balancing equation,
$\widehat S^2=\operatorname{FPdim}(\mathcal E)C$, the full modular
relation, and the requirement that the product central charge of
$\mathcal I\boxtimes\mathcal E$ is one.

The record is
\filename{research/round5/quartic_pair/R6_4_91_center_reconstruction.json}.
The reconstruction and optimizer-free replay are implemented in
\filename{software/round5_quartic_pair.py}. They recompute every packet
cover, orientation, duality, affine projection, normality test, fusion
completion, based-ring orbit, character labelling and twist solution
from the frozen full-tensor frontier. Replay also checks the inherited
signed reduction and its exhaustive profiles when requested.
\end{proof}

The surviving centralizer has no nontrivial invertible object. Its formal
codegree polynomial is
\[
 (t-8)^2(t^2-128t+2048)^2(t^2-64t+512)^2
 (t^4-2816t^3+151552t^2-1572864t+4194304).
\]
For the representative with $z=\zeta_{16}^5$, in the certificate's
unit-first order, the twist exponents are
\[
 (0,10,12,8,14,14,6,6,4,2,8,0,11,11).
\]
The other representative negates these exponents modulo sixteen.
An associator, compatible braiding and the requisite induction algebra
are not supplied by fusion coefficients and modular matrices alone.
Thus the theorem eliminates one inherited profile and completely fixes
the remaining fusion and modular data, while leaving the categorical
realizability question explicit.

\subsection{Full regular-action completion for $R_{6,4,22}$}
\label{subsec:frontier-22-completion}

The remaining quartic ring $R_{6,4,22}$ is treated by the same
restriction--word Gram and regular-action criteria, without a centralizer
factorization. The inherited center ranks are $30$ and $36$.

\begin{proposition}[The rank-$30$ branch]
\label{prop:22-rank30}
The center-rank-$30$ induction profile of $R_{6,4,22}$ is impossible.
\end{proposition}
\begin{proof}
Apply the quadratic and cubic word Gram criterion and couple its coefficient
counts to the projected twist types by
Proposition~\ref{prop:projected-product-balancing}. For each of the two
inherited character labellings the simultaneous system has $2950$
equations and $4132$ nonnegative variables. Each has an exact rational
separating vector. The certificate
\filename{research/round5/fermion/R6_4_22_certificate.json} and
\filename{software/round5_22.py} reconstruct the full domains and check
both separating vectors over the rationals. No integrality or numerical
infeasibility assertion is needed for these two exclusions.
\end{proof}

\begin{theorem}[Complete necessary central data for $R_{6,4,22}$]
\label{thm:22-central-completion}
If $R_{6,4,22}$ has a pseudo-unitary categorification, its center has rank
$36$. Its complete central fusion ring, with the forgetful restriction
map, has one based-ring isomorphism class among all the inherited
possibilities. Each of the two conjugate low character labellings has
exactly one compatible modular datum. The two modular data are complex
conjugate and have twist order $20$.
\end{theorem}
\begin{proof}
For the rank-$36$ profile, retain all $50$ quadratic and cubic words in
the five nonunit low central simples. The complete conditional Gram
ellipsoids of Theorem~\ref{thm:joint-restriction-word-gram} give $155$
possible pairs consisting of a restriction column and its word
coefficient vector. Their simultaneous count system has $902$ equations.
Exact bounded affine enumeration has rank $88$, visits $136389$ nodes,
and gives precisely five nonnegative integral count vectors. The two
low labellings give identical word data, which is verified directly;
thus this count enumeration covers both.

Expand each count vector into the individual central simples. Duality
must dualize its restriction column and permute its word coordinates by
the known low duality. In each of the five patterns this leaves two
possible involutions, corresponding to the two indistinguishable copies
of one column. Apply Proposition~\ref{prop:word-action-completion} to
multiplication by $A_1$. In matrix row convention, its entries satisfy
\[
 (N_{A_1})_{jk}=(N_{A_1})_{k^*,j^*},\qquad
 U_wN_{A_1}=U_{wA_1},
\]
whenever both words have been retained. Tensoriality fixes their
restrictions, and the prescribed eigenvalues fix the full characteristic
polynomial. All variables have finite integral bounds from restriction
products. Exact interval propagation and affine enumeration reject both
dualities in four patterns, and one duality in the fifth pattern.

The last case has $1184$ affine matrices after the prescribed trace is
imposed. Exactly two have the required characteristic polynomial.
Adjoining the actions of $A_2,A_3,A_4$ successively, with commutation
against all previously fixed actions, gives one continuation of each;
$N_{A_5}=N_{A_4}^{\mathsf T}$. The resulting two complete low-action
systems differ only by swapping the indistinguishable central simples.
This equivalence is checked on every matrix entry.

There is a unique central simple whose restriction is $b_1$. Its
multiplication matrix is now determined by restriction, Frobenius
reciprocity and commutation with the fixed low actions: the exact
linear system has $465$ variables, rank $465$, and one nonnegative
integral solution. The central invertible and the other simple actions
are then completed in the same way. Each step has one solution.
The finished tensor is checked directly for the unit, duality,
commutativity, nonnegative integral coefficients and all associativity
identities. Hence the procedure has completed the entire central ring,
not only its products with low simples.

The exact character table of this ring is computed over
$\QQ(\zeta_{20})$ from a separating multiplication matrix. Formal
codegrees, symmetry and the prescribed low character eigenvalues leave
one $S$-labelling for each inherited labelling. Set the six low twists
equal to one. Balancing with these known-twist objects is a linear
system of rank $36$ in the $36$ unknown twists. It has a unique solution
for each labelling; the solutions are complex conjugate and belong to
$\mu_{20}$. The unique nonunit central invertible is forced to have twist $-1$.
Every balancing equation, the full
modular relation and central charge one are checked exactly.

The complete reconstruction is
\filename{software/round5_regular_completion.py}, using the generic
regular-action engine \filename{software/round5_low_actions.py}.
The record is
\filename{research/round5/quartic_pair/R6_4_22_center_reconstruction.json}.
Replay rebuilds the conditional Gram domains, all five count vectors,
all ten duality cases, every surviving matrix, the based-ring
identifications and both modular completions from the frozen full-tensor
frontier. It uses no numerical optimizer and remains valid with Python
assertions disabled.
\end{proof}

The two quartic rings have therefore reached complete necessary central
fusion and modular data. This is stronger than an integral induction or
Gram witness, but it still leaves categorical realization: a compatible
associator and braiding, together with the required canonical induction
algebra, have not been constructed. The current records retain this
remaining question explicitly.

\subsubsection{Finite checks after modular completion}

The completed matrices allow the same rotation criterion to be applied
with every central simple as a source, rather than only the central
unit. Write $S,T$ for the normalized, anomaly-one center matrices and
$W_{Zi}=[F(Z):X_i]$. The generalized indicator vector is
\[
 \bigl(\nu^{Z}_{m,1}(X_i)\bigr)_{Z,i}=ST^mW.
\]
The corresponding rotation on
$\operatorname{Hom}(F(Z),X_i^{\otimes n})$ has $n$th power
$\theta_Z^{-1}$. By \cite[Propositions~3.3--3.7]{BJT}, for
$g=\gcd(n,k)$ its normalized power trace is
\[
 \theta_Z^{k/n}\nu^Z_{n,k}(X_i)
 =\alpha_{k/g,n/g}\!\left(
     \theta_Z^{g/n}\nu^Z_{n/g,1}(X_i^{\otimes g})\right),
\]
where $\alpha_{a,m}(\zeta_m)=\zeta_m^a$. The expression to which this
Galois automorphism is applied must lie in $\QQ(\zeta_{n/g})$.
Exact Fourier inversion must give nonnegative integral eigenspace
dimensions whose sum is
$\dim\operatorname{Hom}(F(Z),X_i^{\otimes n})$.

For both conjugate modular data of each quartic ring, these field
membership and Fourier conditions hold for every original simple,
every central simple, and every $2\le n\le16$. The full center also
satisfies signed-permutation Galois covariance and the twist-square
compatibility. Finally,
\[
 \tau_n=\sum_Z d_Z^2\theta_Z^n>0
 \qquad (\gcd(n,\operatorname{ord}T)=1),
\]
for all eight coprime residue classes of each exponent. Thus every
coprime higher central charge is one, as required of a Drinfeld center
by \cite[Theorem~4.4]{NSWGauss}. Positivity is checked by exact real-root
isolation in the real quartic field. The equality
$\tau_n=\operatorname{FPdim}(\mathcal C)$ is not assumed for general
coprime $n$.

These additional surviving conditions are reconstructed in
\filename{software/round5_modular_checks.py}; its record is
\filename{research/round5/quartic_pair/modular_checks.json}.
They strengthen the verified necessary data and do not resolve the
remaining categorical realizability problem.


The saturated dimension-field characters of both quartic rings have order
four and are nonreal. Their conductors are respectively $20$ and $16$.
The real-character specialization of
Proposition~\ref{prop:saturated-modular-moments} is therefore not applied
to them. Instead, the exact audit
\filename{research/round5/quartic_pair/conductor_intertwiner.py}
constructs an equivariant quotient of
$\operatorname{Ind}_B^{\mathrm{SL}_2(\mathbb Z/f\mathbb Z)}\chi$
onto the cyclic span of $w_\chi$ in each completed modular datum.
It checks every $S,T$ generator edge and
$P_U S w_\chi=\lambda w_{\bar\chi}$, where $P_U$ projects onto
$T$-fixed vectors. The induced and cyclic dimensions are $(36,18)$
for $R_{6,4,22}$ and $(24,24)$ for $R_{6,4,91}$, for both saturated
characters and both modular variants. Complete reducibility and polar
decomposition give an isometric equivariant section of each quotient.
Thus every valid universal conductor moment or spectral identity already
holds on these cyclic spans, and adds no obstruction to the completed data.

The full central projection of
Proposition~\ref{prop:modular-subcategory-projection} also survives.
For the pointed rank-two subcategory of $R_{6,4,22}$ and the Ising
subcategory of $R_{6,4,91}$, twist and forgetful compatibility leave
respectively five and seven integral splitting variables in $B$.
The exact equations $S_{\mathcal D}B=BS_{\mathcal C}$ have ranks five
and seven and force one nonnegative integral matrix in each conjugate
modular variant. Every $S,T$, forgetful and unit equation is rechecked
in \filename{research/round5/quartic_pair/subcategory_projection.py}.
Thus the full projection condition adds no quartic obstruction.
